\section{Evaluation}

\subsection{Datasets}
\label{sec:datasets}

Dataset selection for MACE evaluation proceeds in two stages. The candidate pool is drawn from \citet{lv2025multimodal}, a taxonomy of eighteen multimodal fake news datasets classified by modality availability, classification structure, and source platform.

In the first stage screening, candidate datasets must include: (1) a text modality; (2) a network or propagation modality to support graph-based social context modelling; and (3) a binary or near-binary authenticity classification structure compatible with the MACE output layer. The datasets meeting these criteria are summarised in Appendix~\ref{sec:computeexperiment}, Table~\ref{tab:stage1_candidates}.

In the second stage (Section~\ref{sec:dataset_adaptation}), a detailed adaptation and feasibility assessment determines suitability for Stage 3 experiments; this assessment directly addresses RQ3.

\subsection{Dataset Adaptation}
\label{sec:dataset_adaptation}

All Stage~2 candidate datasets will undergo adaptation to verify suitability for both branches of the MACE architecture; missing or incomplete modalities degrade feature extraction and fusion~\citep{Baltrusaitis2019}. Candidates are evaluated against the following criteria.

\begin{itemize}

    \item \textbf{Text, Image Availability}. Text fields are assessed for availability and completeness; the proportion of records with accessible images will be evaluated to confirm sufficient records for representation learning.

    \item \textbf{Social-Context Availability}. MACE retains the graph-based social-context branches of DANES and GETAE~\citep{truica2024danes, GETAE2025, Song2021}, requiring candidate datasets to provide sufficient user-level metadata, interaction logs, propagation traces, and network connectivity for graph construction.

    \item \textbf{Label Mapping and Class Distribution}. Source labels will be mapped to a unified binary classification structure, and class distribution and imbalance measured across candidates~\citep{Xue2021, mumindataset}.

    \item \textbf{Language and Domain Compatibility}. Non-English datasets (Weibo, CHECKED) require assessment of BERT encoder compatibility. Domain-specific datasets (CoAID, ReCOVery) will be assessed for whether their topical focus on COVID-19 limits generalisability.

\end{itemize}

Stage~2 will produce subset statistics for each dataset, distinguishing the published dataset from the usable subset and confirming the adapted data can support the controlled ablation conditions required by RQ1 and RQ2.

\subsection{Benchmarking}

MACE results will be compared against baselines across three categories:

\begin{itemize}
    \item \textbf{Unimodal with social-context modelling}, consisting of DANES~\citep{truica2024danes} and GETAE~\citep{GETAE2025}, the primary baselines for context-aware text-based detection.

    \item \textbf{Multimodal without social-context modelling}, consisting of HMCAN~\citep{Qian2021}, evaluated in isolation as a comparator during ablation experiments.

    \item \textbf{Multimodal with social-context modelling}, consisting of DPSG~\citep{jing2025} and the multimodal fusion framework of \citet{Dellys2025}.
\end{itemize}

\subsection{Evaluation Metrics}

Stage~3 evaluates the MACE framework across the experimental conditions defined in Section~4, using standard classification metrics consistent with DANES~\citep{truica2024danes}, GETAE~\citep{GETAE2025}, and HMCAN~\citep{Qian2021, su_wan_liu_huang_2020}:

\begin{itemize}
    \item \textbf{Accuracy}, measures the number of correct classifications divided by the total number of classifications.

    \item \textbf{Recall}, measures the number of inauthentic posts correctly classified out of all actual inauthentic posts (true positive rate).

    \item \textbf{Precision}, measures the number of inauthentic posts correctly classified out of all posts classified as inauthentic by the model.

    \item \textbf{F1 Score}, the harmonic mean of recall and precision, balancing false positives and false negatives equally.
\end{itemize}

F1 score will be used as the primary comparison metric across models and datasets, as it captures the balance of detection capability (recall) and reliability (precision). Class imbalance is a known issue in misinformation datasets~\citep{mumindataset}; macro-averaged metrics are therefore reported to ensure balanced consideration of both classes.

Statistical significance testing will be conducted to confirm that observed performance differences are not due to random variation~\citep{stat_test}. McNemar's test~\citep{McNemar_1947} will assess pairwise model comparisons; a paired t-test will evaluate fold-level differences under cross-validation. A significance level of $\alpha = 0.05$ will be used, with performance reported as mean and standard deviation across runs.

\subsection{Computational Efficiency Metrics}

In addition to classification performance, computational efficiency will be evaluated to determine the practical cost of introducing multimodal content processing into the DANES and GETAE frameworks. The following dimensions are reported following \citet{Justus2018}.

\begin{itemize}
    \item \textbf{Training Time}, measured as total wall-clock time per experimental condition, decomposed by batch time and total iterations per epoch.

    \item \textbf{Inference Time}, measured as average prediction time per sample during testing.

    \item \textbf{Peak GPU Memory Usage}, measured as peak memory consumed during model training and inference.

    \item \textbf{Peak System RAM Usage}, measured during graph construction and model execution.

    \item \textbf{Storage Requirements}, covering model checkpoints, cached embeddings, and processed graph structural data.

\end{itemize}

These metrics directly support RQ1 by quantifying the cost of content-branch substitution, and are logged using the evaluation template in Appendix~\ref{sec:computeexperiment}, Table~\ref{tab:efficiency_matrix}.
